% Page 044 — page imprimée 40
% Contenu seul : le préambule, l'en-tête et la pagination sont gérés par meyrueis.tex

temps-là le tourisme constituait un luxe réservé à quelques privilégiés et que les
vacances étaient surtout le fait d’une élite et non pas généralisées comme aujourd’hui.
Les problèmes avaient donc peu de rapport avec ceux rencontrés au Pompidou, sauf
qu’il fallait, à travers un complexe différent en ces deux centres lozériens, faire
entendre le même message de paix, d’amour, de lumière et de justice: le message de
l’Evangile.[\ldots]

\subsection*{Juillet 1932 : Installation à Meyrueis}

Ce mois de juillet fut, cette année là, particulièrement pluvieux. C’est sous les
parapluies et dans la boue que s’effectuèrent nos premiers contacts avec la
population. Notre adaptation ne fut certainement pas facilitée par la clémence du
temps ! J’avais peine à croire que le dépassement de la ligne de partage des eaux qui
séparait le Pompidou de Meyrueis pût amener de tels contrastes de climat. La
Cévenne elle même différait profondément. L’absence de châtaigniers ici était totale
(sauf dans les deux vallées schisteuses et granitiques du Bétuzon et de la Brèze, voire au dessus du
château d’Ayres). Tourné vers l’Atlantique, le paysage se montrait humide, plus
verdoyant, moins abrupt avec cependant, des entailles profondes à la séparation des
Causses, préfigurant déjà les Gorges du Tarn. Bref un changement plus accentué que
prévu.

Cette première impression plutôt fâcheuse fut compensée par notre prise de
possession d’un presbytère vraiment confortable pour l’époque.
Mon prédécesseur le pasteur Albert Chaudier, l’avait laissé en parfait état. Il avait
même voulu, en guise de bienvenue, que les pièces fussent retapissées à neuf pour
nous. Cela nous toucha beaucoup. Quel progrès sur ce que nous quittions ! Eau
courante, cheminées dans toutes les pièces au rez-de-chaussée et au premier,
harmonie dans la disposition de l’ensemble. La cuisine, la salle de séjour et le bureau
donnaient par plusieurs sorties sur un agréable jardin, hautement clôturé, que
remplissaient deux immenses poiriers chargés de fruits. Un conseiller presbytéral
facétieux, me disait « Il conviendrait d’appeler votre demeure « la villa des poires ».
Le moment n’était pas encore venu de posséder une salle de bains installée ni de
toilettes intérieures. Mais ces dernières se trouvaient à proximité et tout pouvait se
passer dans la plus normale discrétion. !

Très conscients de cette situation améliorée, quelle fut notre première préoccupation ?
Aussi inattendue que soit la réponse, la vérité m’oblige à la donner : mettre en route
un enfant.

[ La grossesse fut délicate et le premier-né vit le jour à la Maison de Santé de Nîmes. On l’appela
André, il devint pasteur... comme son père, puis professeur à la faculté de Théologie de Montpellier
écrivit de nombreux livres... ]

Mon ministère fut renouvelé par le bonheur de cette naissance... Je visitai d’arrache-
pied. J’organisai aussi parfaitement que possible mon enseignement auprès des
enfants et des jeunes, écoles du dimanche, du jeudi, catéchismes, Union Chrétienne.
Je montai une chorale et animai des réunions le soir. Je suppléai ainsi à l’absence
momentanée de collaboration de ma femme, absorbée, au début tout au moins, par
son bébé. Mais elle reprit assez vite sa place. Nous avions pu nous faire aider par la
fille d’un conseiller presbytéral, Jeanne Poujol ( fille des Poujol, fermiers de la Fabrique )
